% Generated draft: requires mathematical and source review.
\documentclass[11pt]{article}
\usepackage[margin=25mm]{geometry}
\usepackage{fontspec}
\setmainfont{Latin Modern Roman}
\usepackage{amsmath,amssymb,mathtools,unicode-math}
\setmathfont{Latin Modern Math}
\usepackage{enumitem,needspace,xurl,hyperref}
\setlength{\parindent}{0pt}
\setlength{\parskip}{5pt}
\setlength{\emergencystretch}{3em}
\newcommand{\N}{\mathbb{N}}
\newcommand{\Z}{\mathbb{Z}}
\newcommand{\Q}{\mathbb{Q}}
\newcommand{\R}{\mathbb{R}}
\newcommand{\C}{\mathbb{C}}
\newcommand{\F}{\mathbb{F}}
\newcommand{\A}{\mathbb{A}}
\newcommand{\PP}{\mathbb{P}}
\newcommand{\E}{\mathbb{E}}
\newcommand{\eps}{\varepsilon}
\newcommand{\dd}{\,\mathrm{d}}
\begin{document}
\section*{UnsolvedMath ID: 11000235}
\textbf{Source problem code:} catalogue entry
\subsection{Definitions and mathematical statement}
A pseudo-Anosov mapping class preserves a pair of transverse singular measured foliations and scales their transverse measures by reciprocal factors \(\lambda\) and \(\lambda^{-1}\), with \(\lambda>1\). Problem 4.1 asks whether there exist closed surfaces \(S,S'\), both of genus at least \(2\), and an injective homomorphism
\[
 \rho:\pi_1(S')\hookrightarrow\operatorname{MCG}(S)
\]
satisfying
\[
 \forall\gamma\in\pi_1(S')\setminus\{1\}:\quad
 \rho(\gamma)\text{ is pseudo-Anosov}.
\]
The standard orientable-surface and mapping-class conventions should be checked against the source. The quantifier asks for some pair of surfaces, not every pair or a fixed prescribed genus, and it requires every nonidentity element, not merely every nontrivial generator or simple curve.

The source reports a near-example by Leininger--Reid [LR04], where all closed-curve classes on \(S'\) give pseudo-Anosov elements except one simple closed curve and its iterates. Free homotopy means the exception is understood through the corresponding conjugacy/power classes, with exact conventions requiring the reference. The source then introduces a refinement that is not included in the extract.

The supplied research summary mentions existence of surface subgroups in mapping class groups, names Kahn--Markovic and unspecified earlier constructions, and labels this partial progress. It provides no complete citation or construction verifying the all-nontrivial-elements condition. Existence of an arbitrary surface subgroup is a weaker property and does not answer the stated purely pseudo-Anosov question. The summary's broad affirmative wording is therefore retained only as unverified background, not adopted as a resolution.
\subsection{Short English statement}
Does some closed surface group embed in the mapping class group of another closed surface so that every nonidentity element is pseudo-Anosov? Both genera must be at least two. The supplied near-example has an exceptional curve and powers; general surface-subgroup existence does not settle this stronger condition.
\subsection{Sources}
\begin{itemize}
\item Local UnsolvedMath registry, permanent ID \(11000235\), exact source code \texttt{catalogue entry}, Problem 4.1, PDF page 272. No individual proposer or proposed year is supplied. The extract includes a page footer naming L. Mosher and an introduction to a missing refinement.
\item Farb, edited volume \emph{Problems on Mapping Class Groups and Related Topics}: \url{https://www.math.uchicago.edu/~farb/papers/mcgbook.pdf}. The volume has 329 explicitly typeset numbered items across chapters, repeated local numbering distinguished by PDF page.
\item Leininger--Reid [LR04] for the source-reported near-example; full entry absent. The research summary's Kahn--Markovic and earlier-construction mentions have no supplied precise citation and no verified relation to the purely pseudo-Anosov target.
\item Extraction pdf-font-structure, accessed 2026-07-29. Source evidence NEEDS\_REVIEW presents an open problem and requires a current primary-literature audit. Registry status open and research classification OPEN-TRIAGE coexist with a summary labeled PARTIAL about the weaker surface-subgroup property.
\item Dataset attribution: Ulam AI / UnsolvedMath Contributors, CC BY 4.0. Source rights NEEDS\_REVIEW: publicly accessible author/publisher copy, redistribution terms requiring review.
\end{itemize}
\end{document}
